一量子比特的基态为 \(|0\rangle\)、激发态为 \(|1\rangle\)，哈密顿量为 \(H=\hbar\omega|1\rangle\langle1|\)，其中 \(\omega>0\)。零温环境起初处于真空 \(|0_E\rangle\)，在无外部驱动的马尔可夫近似下，比特自发辐射的速率为 \(\Gamma>0\)。令 \(\eta=e^{-\Gamma t}\in(0,1]\)，总演化在单激发子空间内守恒激发数，且
\(U|0,0_E\rangle=|0,0_E\rangle\)，
\(U|1,0_E\rangle=\sqrt{\eta}|1,0_E\rangle+\sqrt{1-\eta}|0,1_E\rangle\)。
(1) 将上述等距映射扩展为单激发子空间的酉变换，迹掉环境推导振幅阻尼通道的两个 Kraus 算符，验证完备性，并求一般初态密度矩阵及 Bloch 向量的演化。
(2) 从通道的连续时间形式推导零温 Lindblad 主方程，求相干元和激发态布居的精确衰减，证明通道半群性质并核对比特能量的耗散率。
(3) 取参考比特 \(A\) 与系统 \(B\) 初态为 \(|\Phi^+\rangle=(|00\rangle+|11\rangle)/\sqrt2\)，只对 \(B\) 施加通道。写出末态 \(4\times4\) 密度矩阵，计算对 \(B\) 的部分转置谱及负性 \(\mathcal N=(\|\rho^{T_B}\|_1-1)/2\)，判断有限时间是否发生纠缠突然死亡。
(4) 对任意纯态 \(|\psi\rangle=\alpha|0\rangle+\beta|1\rangle\)，求输入输出态保真度 \(F_\psi\)。对 Bloch 球 Haar 均匀平均求 \(F_{\rm av}\)，再独立用最大混合态的纠缠保真度 \(F_e\) 验证 \(F_{\rm av}=(2F_e+1)/3\)，检查 \(t=0\) 和 \(t\to\infty\) 极限。
